\origpage{338}{N. Robertson, P.D. Seymour / Journal of Combinatorial Theory, Series B 92 (2004) 325--357}
\hspace*{12pt}(1) \textit{For each $(A, B) \in \mathcal{M}(\mathcal{T}, W', \theta)$, there is no $(C, D) \in \mathcal{L}^-$ with $A \subseteq C$ and $D \subseteq B$.}

\textit{Subproof:} Let $(A, B) \in \mathcal{M}(\mathcal{T}, W', \theta)$ and suppose that there is such a $(C, D)$. Since $(A, B) \in \mathcal{M}(\mathcal{T}, W', \theta)$ there is a tangle $\mathcal{T}'$ such that $(A, B)$ is the $(\mathcal{T}, \mathcal{T}')$-distinction, and there exists $z \in W' \setminus Z(B)$. Now from the minimality of $W'$, $\mathcal{L}$ does not $\theta$-isolate $\mathcal{T}$ modulo $W' \setminus \{z\}$, and so there exists $(C', D') \in \mathcal{L}^-$ and a tangle $\mathcal{T}''$ in $G$ of order $\ge \theta$ with $(D', C') \in \mathcal{T}''$ with the property that $W' \setminus \{z\} \subseteq Z(B')$ and not both $A' \subseteq C'$ and $D' \subseteq B'$, where $(A', B')$ is the $(\mathcal{T}, \mathcal{T}'')$-distinction. Since $\mathcal{L}$ does $\theta$-isolate $\mathcal{T}$ modulo $W'$, it follows that $W' \nsubseteq Z(B')$ and so $z \notin Z(B')$. Hence $B \cup B' \ne G$. Moreover, since $(A, B) \in \mathcal{M}(\mathcal{T}, W', \theta)$, it follows that not both $A \subseteq A'$ and $B' \subseteq B$, from the third condition in the definition of $\mathcal{M}(\mathcal{T}, W', \theta)$. From 5.1, $A' \subseteq A$ and $B \subseteq B'$. Now $A \subseteq C$ and $D \subseteq B$, and so $A' \subseteq C$ and $D \subseteq B'$; and hence $(C, D) \ne (C', D')$, since not both $A' \subseteq C'$ and $D' \subseteq B'$. Moreover, $(B', A') \in \mathcal{T}''$, and $A' \subseteq C$ and $D \subseteq B'$, and so $(D, C) \in \mathcal{T}''$ since $(D, C)$ has order $< \theta$ and $\mathcal{T}''$ has order $\ge \theta$. Since $\mathcal{L}$ is a rooted location and $(C, D), (C', D') \in \mathcal{L}^-$ it follows that $D \cup D' = G^-$. But $(D, C), (D', C') \in \mathcal{T}''$ contrary to the second axiom for tangles. This proves (1).

Let $X = \bigcup(V(A \cap B) : (A, B) \in \mathcal{M}(\mathcal{T}, W', \theta))$. Since $|\mathcal{M}(\mathcal{T}, W', \theta)| \le |W'| \le |W| \le \theta$ by 5.2, it follows that $|X| \le \theta(\theta - 1)$. For each $C \in \mathcal{L}$, let $f(C)$ be a rooted hypergraph with $f(C)^- = C^-$ and $\bar{\pi}(f(C)) = \bar{\pi}(C) \cup (X \cap V(C))$, taking $f(C) = C$ if $X \cap V(C) \subseteq \bar{\pi}(C)$. Let $\mathcal{L}' = \{f(C) : C \in \mathcal{L}\}$. Then $\mathcal{L}'$ is a fine rooted location, and $\mathcal{L}'$ has order at most $\theta(\theta - 1)$ more than the order of $\mathcal{L}$, and hence at most $\theta^2 - 1$. We observe

\hspace*{12pt}(2) \textit{For each $(C', D') \in \mathcal{L}'^-$ there exists $(C, D) \in \mathcal{L}^-$ with $C = C'$ and $D \subseteq D'$; and $E(D') = E(D)$, and $V(D') \setminus V(D) = X \cap (V(C) \setminus \bar{\pi}(C))$.}

Since $|X| \le \theta(\theta - 1)$ and each $x \in X$ belongs to $V(C) \setminus \bar{\pi}(C)$ for at most one $C \in \mathcal{L}$, we see that $\mathcal{F}'$ contains a heart of $(P, \mathcal{L}')$, from the definition of $\mathcal{F}'$. Since $\mathcal{T}$ has order $\ge \theta^2$ and $\mathcal{L}^- \subseteq \mathcal{T}$ it follows from (2) that $\mathcal{L}'^- \subseteq \mathcal{T}$. To verify that $\mathcal{L}'$ is $W'$-suitable, let $\mathcal{T}'$ be a tangle of order $\ge \theta^2$, let $(C', D') \in \mathcal{L}'^-$ with $(D', C') \in \mathcal{T}'$, and let $(A, B)$ be the $(\mathcal{T}, \mathcal{T}')$-distinction. We may assume that:

\hspace*{12pt}(3) \textit{There is no $(A^*, B^*) \in \mathcal{M}(\mathcal{T}, W', \theta)$ such that $A \subseteq A^*$ and $B^* \subseteq B$.}

We must therefore show that $A \subseteq C'$ and $D' \subseteq B$. Choose $(C, D)$ as in (2). Then $(A, B)$ has order at most that of $(C, D)$, and hence $< \theta$. If $W' \nsubseteq Z(B)$, then from the definition of $\mathcal{M}(\mathcal{T}, W', \theta)$, there exists some $(A^*, B^*) \in \mathcal{M}(\mathcal{T}, W', \theta)$ violating (3); so $W' \subseteq Z(B)$. Since $\mathcal{L}$ $\theta$-isolates $\mathcal{T}$ modulo $W'$ and $(D, C) \in \mathcal{T}'$, it follows that $A \subseteq C$ and $D \subseteq B$. Since $C = C'$ it remains to show that $D' \subseteq B$. Let $v \in V(D') \setminus V(D)$. Then $v \in X$, and so $v \in V(A^* \cap B^*)$ for some $(A^*, B^*) \in \mathcal{M}(\mathcal{T}, W', \theta)$. By 5.1, (3) and the third condition in the definition of $\mathcal{M}(\mathcal{T}, W', \theta)$, it follows that $A \subseteq B^*$ and $A^* \subseteq B$; and in particular $v \in V(B)$. Consequently $V(D') \setminus V(D) \subseteq V(B)$; and since $E(D') = E(D)$ and $D \subseteq B$, it follows that $D' \subseteq B$ as required. This proves that $\mathcal{L}'$ is $W'$-suitable. The final statement holds because of (1) and the definition of $\mathcal{L}'$. $\square$

If $x, y$ are vertices of a graph $H$, we say they are \textit{connected in $H$} if they belong to the same connected component of $H$.
